\section{Introduction}

A retrieval-augmented generation (RAG) service is a pipeline: a retriever finds passages, glue code
assembles them into a prompt, a serving runtime runs the model, and an index holds the documents.
Evaluation frameworks for RAG score the retrieved context and the final
answer~\cite{es2024ragas,saadfalcon2024ares,ru2024ragchecker}. In production, however, evidence can
be lost \emph{between} those two points. A government help-desk assistant lost its system prompt and
its top-ranked documents for weeks because the serving runtime silently kept only the tail of a
prompt that exceeded its default context window~\cite{musetoiu2026ragal}; users of a popular RAG
platform report the same problem when the model server ignores the configured
window~\cite{ragflow_ollama_ctx}. Template code that truncates passages, an index that still holds
an outdated version of an updated document~\cite{ouyang2025hoh}, and a query encoder that no longer
matches the encoder the index was built with are similar failures: the retriever may work, the model
may work, and the answer is still wrong.

When answers look wrong, a common runtime response is generic: retry, or ask the model whether its
answer is supported and regenerate if not. Such repairs spend calls on the generator even when the
generator is not the component that failed, and LLMs are known to be unreliable at verifying their
own outputs without external feedback~\cite{huang2024selfcorrect}. This paper asks a practical
question for teams operating RAG services: \emph{when a pipeline component fails, which runtime
repair helps, what does it break, and what does it cost?}

We study this with a fault-injection harness around a small but complete RAG service running
locally. We inject six faults motivated by reported incidents at four pipeline stages (retrieval,
prompt assembly, serving, index content), and compare twelve repair policies under a fixed budget
of two extra model calls. The policies range from generic (retry, self-check) to
\emph{trace-based}: checks computed from what the service already logs, such as whether the rendered
prompt is shorter than the retrieved text or whether the server processed fewer tokens than were
sent. We measure success under faults, damage to answers that were correct, abstention when the
evidence is gone, and cost. We address four questions:
\begin{description}
\item[RQ1] How do generic, blanket and trace-based repairs compare under component faults?
\item[RQ2] Does the verifier or the repair action matter more? We cross \{LLM self-verifier,
MiniCheck grounding classifier~\cite{tang2024minicheck}, text match\} with \{regenerate on the
same passages, repair the context\}.
\item[RQ3] How quickly do trace checks and model-based signals detect a faulty deploy in a
request stream?
\item[RQ4] Can a regression suite built from observed failures catch regressions in the repair
logic itself?
\end{description}

On 150 held-out HotpotQA questions with Llama-3.2-3B, trace-based diagnosis raises success under
faults from 26.8\% to 56.9\% (+30.1 points, 95\% CI [24.8, 35.7]) while breaking 4.7\% of correct
clean answers; trace checks alone account for almost all of the gain. A generic LLM self-check
lowers success to 20.7\% and breaks 24.3\% of correct clean answers. Applying every repair on every
request reaches 46.4\%, breaks 15.9\% of correct clean answers, and makes the serving fault worse
than doing nothing (6.0\% vs.\ 30.0\%). Trace checks detect each global fault a median of 4
requests after a deploy; the model's own verifier needs tens of requests and often never alarms.
A 35-case regression suite, replayed from recorded responses in seconds, fails whenever one of
the diagnostics is disabled.

\paragraph{Contributions.} (1) A fault-injection study of RAG runtime repair that treats prompt
assembly and serving as components that can fail; (2) a budget-matched comparison of twelve
repair policies that accounts for damage and abstention; (3) evidence that cheap, model-free trace
checks recover most of the benefit and detect faulty deploys fastest, while model self-verification
is harmful at this scale; and (4) an open, laptop-scale artifact with record/replay regression tests
for LLM pipelines.
